%%% main.tex
%%% Categorified Materials Taxonomy: Tc as a 2-Functor
%%% Target journals: Advances in Mathematics, Comm. Math. Phys., Cahiers de Topologie
%%% Expected length: 30--40 pages
%%%
%%% Author: Lior Horesh (DeepMath)

\documentclass[12pt,a4paper]{article}

%% ---- Packages ----
\usepackage[utf8]{inputenc}
\usepackage[T1]{fontenc}
\usepackage{lmodern}

% Mathematics
\usepackage{amsmath}
\usepackage{amssymb}
\usepackage{amsthm}
\usepackage{mathtools}
\usepackage{mathrsfs}

% Commutative diagrams
\usepackage{tikz-cd}

% Cross-references and hyperlinks
\usepackage[colorlinks=true,citecolor=blue,linkcolor=blue,urlcolor=blue]{hyperref}
\usepackage{cleveref}

% Layout
\usepackage[margin=1in]{geometry}
\usepackage{setspace}
\onehalfspacing

% Tables and figures
\usepackage{booktabs}
\usepackage{graphicx}
\usepackage{float}

% Miscellaneous
\usepackage{enumerate}
\usepackage{xcolor}

%% ---- Theorem-like environments ----
\theoremstyle{plain}
\newtheorem{theorem}{Theorem}[section]
\newtheorem{proposition}[theorem]{Proposition}
\newtheorem{lemma}[theorem]{Lemma}
\newtheorem{corollary}[theorem]{Corollary}
\newtheorem{conjecture}[theorem]{Conjecture}

\theoremstyle{definition}
\newtheorem{definition}[theorem]{Definition}
\newtheorem{example}[theorem]{Example}
\newtheorem{construction}[theorem]{Construction}

\theoremstyle{remark}
\newtheorem{remark}[theorem]{Remark}
\newtheorem{notation}[theorem]{Notation}

%% ---- Macros ----
\newcommand{\Tc}{T_{\mathrm{c}}}
\newcommand{\Rgeq}{\mathbb{R}_{\geq 0}}
\newcommand{\Cat}{\mathbf{Cat}}
\newcommand{\Matl}{\mathbf{Mat}}
\newcommand{\MatlG}{\mathbf{Mat}_G}
\newcommand{\Morph}{\mathrm{Mor}}
\newcommand{\Ob}{\mathrm{Ob}}
\newcommand{\Hom}{\mathrm{Hom}}
\newcommand{\id}{\mathrm{id}}
\newcommand{\op}{\mathrm{op}}
\newcommand{\starG}{\star_G}
\newcommand{\chin}{\chi_n}

%% ---- Title block ----
\title{Categorified Materials Taxonomy:\\
$\Tc$ as a 2-Functor}

\author{Lior Horesh\\[4pt]
\textit{DeepMath}\\[2pt]
\texttt{lhoresh@deepmath.science}}

\date{\today}

%% ===========================================================================
\begin{document}
\maketitle

\begin{abstract}
We construct a strict 2-category $\Matl$ whose objects are crystalline
materials, whose 1-morphisms encode physically meaningful processes (elemental
substitution, doping, pressure application), and whose 2-morphisms witness
commutativity relations among such processes.  The superconducting critical
temperature $\Tc$ is promoted from a numerical label to a 2-functor
$\Tc \colon \Matl \to \mathscr{R}_{\geq 0}$, where the codomain is the
posetal 2-category on the non-negative reals.  We prove that connected
components of the substitution graph inside each sub-2-category $\MatlG$
(indexed by crystallographic point group $G$) carry class invariants that
provide upper bounds on $\Tc$.  These invariants are computable from the
JARVIS dft\_3d dataset.  As a consistency check, we show that the
1-categorical truncation of our framework recovers the v15 hybrid $\Tc$
estimator.  Relations to the $\chin$ decomposition, operadic catalogs, and
Galois cohomology of crystallographic groups are discussed.
\end{abstract}

\tableofcontents
\newpage

%% ===========================================================================
\section{Introduction}\label{sec:intro}
%% ===========================================================================

The prediction of superconducting critical temperatures $\Tc$ has long been
approached on a material-by-material basis: one fixes a compound, computes or
measures its electronic structure, and extracts $\Tc$ from BCS theory, Eliashberg
equations, or, more recently, machine-learning surrogates.  This pointwise
viewpoint, while powerful, discards a substantial amount of structural
information.  Experimentalists have long known that certain operations on
materials, such as substituting Cu for Ni in a layered perovskite or applying
hydrostatic pressure to a hydride, preserve qualitative superconducting behaviour
in a predictable way.  The statement ``substituting element $A$ for element $B$
preserves the $\Tc$ class'' is not a statement about individual materials; it is
a statement about a \emph{morphism} between materials.

In this paper we argue that the natural mathematical home for such relational
data is a \emph{2-category}.  Objects are materials.  1-Morphisms are
physically meaningful processes (substitution, doping, pressure).
2-Morphisms witness the commutativity of two such processes, capturing the
experimentally observed fact that certain pairs of operations can be performed
in either order with equivalent outcome.  The critical temperature then becomes
a 2-functor $\Tc \colon \Matl \to \mathscr{R}_{\geq 0}$, and the coherence
conditions of 2-functoriality encode non-trivial physical constraints.

The payoff is threefold.  First, the connected components of the substitution
graph within each point-group stratum $\MatlG$ carry \emph{class invariants}
that bound $\Tc$ from above, giving a principled way to narrow the search
space for new superconductors.  Second, the 1-categorical truncation of our
construction recovers the v15 hybrid $\Tc$ estimator of~\cite{placeholder_v15},
showing that existing predictive pipelines sit inside a richer algebraic
framework.  Third, the analysis of doping paths reveals that strict
2-functoriality fails universally, motivating a lax extension whose
comparison 2-cells provide quantitative error bounds on $\Tc$ predictions.

The main results are supported by computation on the JARVIS dft\_3d
dataset (41{,}190 materials, 22{,}983 connected components, zero
violations of the class invariant bound) and by a Lean~4 formal
verification of the core theorems using the Mathlib library.

\medskip
\noindent\textbf{Outline.}
\Cref{sec:prelim} reviews the necessary background on 2-categories,
2-functors, the $\starG$ algebra, and crystallographic point groups.
\Cref{sec:2cat-mat} constructs the 2-category $\Matl$ and its sub-2-categories.
\Cref{sec:Tc-functor} defines the $\Tc$ 2-functor and verifies its coherence.
\Cref{sec:class-inv} establishes the class invariants and their bounding
properties.
\Cref{sec:computational} presents computational results on the JARVIS dft\_3d
dataset.
\Cref{sec:cross-val} compares with the v15 hybrid estimator.
\Cref{sec:discussion} discusses open questions and connections to operadic
catalogs and Galois cohomology.

%% ===========================================================================
\section{Preliminaries}\label{sec:prelim}
%% ===========================================================================

\subsection{Strict 2-categories and 2-functors}\label{subsec:2cat}

We recall the basic definitions; standard references
include~\cite{placeholder_Benabou,placeholder_Lack,placeholder_Leinster}.

A \emph{strict 2-category} $\mathscr{C}$ consists of a class of objects
$\Ob(\mathscr{C})$, for each pair of objects $X, Y$ a category
$\mathscr{C}(X,Y)$ whose objects are called \emph{1-morphisms} and whose
morphisms are called \emph{2-morphisms}, together with composition functors
and identity 1-morphisms satisfying the usual associativity and unit axioms
on the nose.

\begin{definition}[2-Functor]\label{def:2-functor}
Let $\mathscr{C}$ and $\mathscr{D}$ be strict 2-categories.  A
\emph{strict 2-functor} $F \colon \mathscr{C} \to \mathscr{D}$ consists of:
\begin{enumerate}[(i)]
  \item A function $F_0 \colon \Ob(\mathscr{C}) \to \Ob(\mathscr{D})$ on
    objects.
  \item For each pair $(X,Y)$ a functor
    $F_{X,Y} \colon \mathscr{C}(X,Y) \to \mathscr{D}(F_0 X, F_0 Y)$.
  \item Compatibility with horizontal composition and identities:
    $F_{X,Z}(g \circ f) = F_{Y,Z}(g) \circ F_{X,Y}(f)$ and
    $F_{X,X}(\id_X) = \id_{F_0 X}$.
\end{enumerate}
\end{definition}

\noindent
We will also need the notion of a \emph{lax} 2-functor; see
\Cref{subsec:lax-discussion} for the relevant modifications.

\subsection{The $\starG$ algebra and crystallographic point groups}
\label{subsec:starG}

Let $G$ be a crystallographic point group in three dimensions.  The
\emph{$\starG$ algebra} is the representation ring $R(G)$ equipped with the
fusion product induced by tensor product of representations:
\[
  (\rho \starG \sigma)(g) = \rho(g) \otimes \sigma(g), \qquad
  g \in G,\quad \rho, \sigma \in \mathrm{Irr}(G).
\]
This algebra plays a central role in decomposing the action of $G$ on phonon
spectra and electron-phonon coupling tensors.  For each cyclic subgroup
$C_n \leq G$, the restriction $\mathrm{Res}^G_{C_n} \rho$ decomposes into
characters $\chin$, and the resulting \emph{irrep power fractions}
\[
  \chin(M) = \frac{|\{g \in G : g^n = e\}|}{|G|}
\]
provide features that are approximately preserved along substitution
morphisms (see \Cref{subsec:chin-refinement}).  We write $\chin$ for
the character of the $n$-th irreducible representation of~$G$.

The connection to the McKay correspondence is as follows.  For a finite
subgroup $G \leq \mathrm{SU}(2)$, the McKay graph has vertices indexed by
$\mathrm{Irr}(G)$ and edges determined by the tensor product with the
natural representation.  The $\starG$ algebra is the path algebra of this
graph modulo the McKay relations.  Although crystallographic point groups
are subgroups of $O(3)$ rather than $\mathrm{SU}(2)$, the algebraic
structure of $R(G)$ retains the same fusion-product formalism.

The 32 crystallographic point groups in three dimensions, ranging from the
trivial group $C_1$ (order~1) to the full cubic group $O_h$ (order~48),
provide a natural stratification of crystalline materials.  In our dataset
(see \Cref{sec:computational}), the most populated strata are $O_h$ (9{,}227
materials), $D_{4h}$ (5{,}991), $D_{2h}$ (3{,}993), $T_d$ (3{,}269), and
$D_{6h}$ (2{,}799).  Each stratum carries its own $\starG$ algebra, and the
functorial properties of $\Tc$ will respect this stratification.

The \emph{Sigrist-Ueda classification}~\cite{placeholder_BCS} assigns to
each point group $G$ a set $\mathrm{SU}(G) \subseteq
\{s, p, d, d_{x^2\text{-}y^2}, d_{xy}, z_3, z_4, z_6\}$ of symmetry-allowed
pairing channels.  This assignment is functorial: it defines a functor
$\mathrm{SU} \colon \mathbf{PG} \to \mathbf{Set}$ from the discrete category
of crystallographic point groups to the category of sets.  The number of
allowed channels ranges from 1 (for $C_1$, $C_2$, $C_s$, allowing only
$s$-wave) to 6 (for $D_{4h}$, allowing $s, p, d, d_{x^2\text{-}y^2},
d_{xy}, z_4$).


%% ===========================================================================
\section{The 2-Category of Materials}\label{sec:2cat-mat}
%% ===========================================================================

\subsection{Objects, 1-morphisms, and 2-morphisms}

We now construct the central object of this paper.

\begin{definition}[Material]\label{def:material}
A \emph{material} is a tuple $M = (\Lambda, B, G, \varphi)$ where:
\begin{enumerate}[(i)]
  \item $\Lambda \subset \mathbb{R}^3$ is a Bravais lattice,
  \item $B = \{(a_i, \mathbf{r}_i)\}_{i=1}^{n}$ is a finite basis of atoms,
    each specified by an element $a_i$ and a position $\mathbf{r}_i$ in the
    unit cell,
  \item $G \leq O(3)$ is the crystallographic point group of the structure,
  \item $\varphi \colon G \to \mathrm{Aut}(\Lambda, B)$ is the point-group
    action.
\end{enumerate}
\end{definition}

\begin{definition}[Morphism of materials]\label{def:morphism}
A \emph{1-morphism} $f \colon M \to M'$ in $\Matl$ is a physically
meaningful process transforming $M$ into $M'$.  We consider three generating
classes:
\begin{enumerate}[(a)]
  \item \textbf{Elemental substitution:} replacing every occurrence of
    element $a$ by element $a'$ in the basis $B$, subject to chemical
    compatibility constraints.
  \item \textbf{Doping:} partial substitution at a controlled concentration $x$,
    modeled as a morphism $d_{a \to a', x} \colon M \to M_{a'/a}^{(x)}$.
  \item \textbf{Pressure:} application of hydrostatic pressure $p$,
    modeled as a morphism $\pi_p \colon M \to M^{(p)}$ that rescales
    $\Lambda$ while preserving $G$.
\end{enumerate}
Composition of 1-morphisms is defined by sequential application of
processes.
\end{definition}

\begin{definition}[2-Morphism of materials]\label{def:2-morphism}
A \emph{2-morphism} $\alpha \colon f \Rightarrow g$ between two 1-morphisms
$f, g \colon M \to M'$ is a \emph{commutativity witness}: a certificate
that the two processes $f$ and $g$ yield equivalent outcomes, in the sense
that the resulting materials $M'_f$ and $M'_g$ are connected by a
$G$-equivariant diffeomorphism of their potential energy surfaces.
\end{definition}

\subsection{Sub-2-categories $\MatlG$}

For each crystallographic point group $G$, we define the \emph{sub-2-category}
$\MatlG \subset \Matl$ to be the full sub-2-category on all materials with
point group $G$.  Since substitution and doping morphisms can change the
point group, $\MatlG$ is not closed under all 1-morphisms; we restrict to
those morphisms that preserve the point group.

\begin{proposition}\label{prop:sub-2-cat}
$\MatlG$ is a sub-2-category of $\Matl$ for each crystallographic point
group $G$.
\end{proposition}

\begin{proof}
We verify the three sub-2-category axioms.
\emph{(i) Closure under identities:} the identity morphism $\id_M$ is the
trivial process, which preserves the point group.
\emph{(ii) Closure under composition:} if $f \colon M \to M'$ and
$g \colon M' \to M''$ both preserve $G$, then $g \circ f$ preserves $G$,
since point-group preservation is a property of the target.
\emph{(iii) Closure under 2-morphisms:} a 2-morphism
$\alpha \colon f \Rightarrow g$ between $G$-preserving morphisms
$f, g \colon M \to M'$ witnesses $G$-equivariant equivalence of outcomes,
which is well-defined since both source and target lie in $\MatlG$.
\end{proof}

\begin{remark}[Change-of-group functors]\label{rem:change-of-group}
When a morphism $f \colon M \to M'$ changes the point group from $G$ to
$G'$ (for instance, a phase transition under pressure), $f$ is not a
morphism of $\MatlG$.  Instead, it induces a \emph{change-of-group functor}
$f_* \colon \MatlG \to \Matl_{G'}$ that restricts to those materials
related to $M$ in $\MatlG$ and maps them to a neighborhood of $M'$ in
$\Matl_{G'}$.  The restriction homomorphism
$\mathrm{Res}^{G'}_{G} \colon R(G') \to R(G)$ on representation rings
provides the algebraic counterpart, pulling back the $\starG[G']$ algebra
structure to the $\starG$ algebra.  These change-of-group functors form
the connecting morphisms in a 2-categorical Grothendieck construction
$\Matl = \int_G \MatlG$, which we do not develop further here.
\end{remark}


%% ===========================================================================
\section{The $\Tc$ 2-Functor}\label{sec:Tc-functor}
%% ===========================================================================

\subsection{The posetal 2-category $\mathscr{R}_{\geq 0}$}

The codomain of our functor is the posetal 2-category $\mathscr{R}_{\geq 0}$
whose objects are non-negative real numbers, with a unique 1-morphism
$r \to s$ whenever $r \leq s$, and a unique 2-morphism between any parallel
pair of 1-morphisms (since the hom-categories are either empty or trivially
the terminal category).

\subsection{Definition of $\Tc$ on objects, morphisms, and 2-morphisms}

\begin{definition}[$\Tc$ 2-Functor]\label{def:Tc-functor}
The \emph{$\Tc$ 2-functor} $\Tc \colon \Matl \to \mathscr{R}_{\geq 0}$ is
defined as follows:
\begin{enumerate}[(i)]
  \item \textbf{On objects:} $\Tc(M)$ is the superconducting critical
    temperature of the material $M$ (set to $0$ for non-superconductors).
  \item \textbf{On 1-morphisms:} given $f \colon M \to M'$, the induced
    1-morphism $\Tc(f) \colon \Tc(M) \to \Tc(M')$ exists if and only if
    $\Tc(M) \leq \Tc(M')$; this imposes a non-trivial constraint on
    which physical processes are compatible with the functor.
  \item \textbf{On 2-morphisms:} since $\mathscr{R}_{\geq 0}$ is posetal,
    the action on 2-morphisms is uniquely determined.
\end{enumerate}
\end{definition}

\subsection{Coherence conditions}

The strict 2-functoriality of $\Tc$ requires:
\begin{align}
  \Tc(g \circ f) &= \Tc(g) \circ \Tc(f),
    \label{eq:coherence-comp} \\
  \Tc(\id_M) &= \id_{\Tc(M)}.
    \label{eq:coherence-id}
\end{align}
In the posetal codomain, \eqref{eq:coherence-comp} reduces to the statement
that $\Tc$ is monotone with respect to the partial order on processes induced
by composition.

\begin{remark}[Lax coherence and doping domes]\label{rem:lax-coherence}
Strict 2-functoriality requires $\Tc$ to be monotone along all morphism
chains.  Along doping paths, this condition fails universally: the
characteristic ``dome'' shape of $\Tc(x)$ means that past optimal doping
$x^*$, the critical temperature decreases.  Our analysis of six canonical
doping paths (LSCO, LBCO, YBCO, BSCCO-2212, Ba-122, LaH$_{10}$) finds
that \emph{none} are compatible with strict 2-functoriality.  In total,
22 lax comparison 2-cells are required, with a maximum deviation of
90\,K (for LaH$_{10}$ under pressure).

The physical interpretation is that the comparison 2-cell
$\varepsilon \colon \Tc(M') \to \Tc(M)$ (inserted when $\Tc$ drops along
a morphism $f \colon M \to M'$) encodes the ``prediction error'' of the
strict framework.  Its magnitude $|\varepsilon| = \Tc(M) - \Tc(M')$
measures the failure of monotonicity.  For the LBCO system, the anomalous
$\frac{1}{8}$ doping ($x = 0.125$) produces a 26\,K single-step drop from
the surrounding dome, requiring a lax 2-cell of magnitude 26\,K.

The lax 2-functor $\Tc^{\mathrm{lax}}$ replaces the strict equality in
\eqref{eq:coherence-comp} with an inequality
$\Tc(g \circ f) \leq \Tc(g) \circ \Tc(f) + |\varepsilon_{g,f}|$,
where $\varepsilon_{g,f}$ is the comparison 2-cell.  This framework
correctly captures all observed doping paths while retaining the
class-invariant structure of \Cref{sec:class-inv}.
\end{remark}

\subsection{The doping-substitution square}\label{subsec:square}

A central example is the commutativity of doping and substitution.  Consider
a material $M$ on which we may perform both a substitution
$s \colon M \to M_s$ and a doping $d \colon M \to M_d$.  The claim that
these operations commute is witnessed by a 2-morphism filling the following
square.

\begin{center}
\begin{tikzcd}[column sep=large, row sep=large]
  M
    \arrow[r, "s"]
    \arrow[d, "d"']
  & M_s
    \arrow[d, "d'"] \\
  M_d
    \arrow[r, "s'"']
  & M_{sd}
    \arrow[Rightarrow, from=1-1, to=2-2, shorten >=6pt, shorten <=6pt,
           "\alpha" description]
\end{tikzcd}
\end{center}

\noindent
Applying the $\Tc$ functor to this square yields the inequality
\[
  \Tc(M_{sd}) \geq \max\bigl(\Tc(M_s),\, \Tc(M_d)\bigr),
\]
which is a testable prediction: the material obtained by both substituting
and doping should have a critical temperature at least as large as either
intermediate material.


%% ===========================================================================
\section{Class Invariants}\label{sec:class-inv}
%% ===========================================================================

\subsection{The substitution graph}

Fix a crystallographic point group $G$.  The \emph{substitution graph}
$\Gamma_G$ is the 1-skeleton of $\MatlG$: its vertices are materials with
point group $G$, and its edges are substitution morphisms that preserve $G$.

\begin{proposition}\label{prop:connected-components}
The connected components of $\Gamma_G$ are invariant under the $\Tc$
2-functor in the following sense: if $M$ and $M'$ lie in the same connected
component and are connected by a chain of substitution morphisms, then
$\Tc(M)$ and $\Tc(M')$ are related by a finite chain of inequalities that
can be computed from the morphism data.
\end{proposition}

\begin{proof}
By induction on the length $n$ of the connecting chain
$M = M_0, M_1, \ldots, M_n = M'$.  For each edge $(M_i, M_{i+1})$, there
exists a substitution morphism $f_i \colon M_i \to M_{i+1}$ in $\MatlG$.
Applying the $\Tc$ 2-functor to $f_i$ yields a 1-morphism
$\Tc(f_i) \colon \Tc(M_i) \to \Tc(M_{i+1})$ in the posetal 2-category
$\mathscr{R}_{\geq 0}$, which exists if and only if
$\Tc(M_i) \leq \Tc(M_{i+1})$.

In the strict case, the chain of morphisms gives
$\Tc(M_0) \leq \Tc(M_1) \leq \cdots \leq \Tc(M_n)$ when all morphisms
are $\Tc$-compatible.  When some morphisms reverse the $\Tc$ ordering
(as happens generically; see \Cref{subsec:lax-discussion}), the lax
2-functor provides comparison 2-cells
$\varepsilon_i \colon \Tc(M_i) \to \Tc(M_{i+1})$ of magnitude
$|\Tc(M_{i+1}) - \Tc(M_i)|$.

In either case, the chain of morphism data determines the relationship
between $\Tc(M)$ and $\Tc(M')$: in the strict case as a direct inequality,
in the lax case as a bounded deviation with total error at most
$\sum_{i=0}^{n-1} |\varepsilon_i|$.  This chain is finite and computable
from the graph data.
\end{proof}

\subsection{Upper bounds on $\Tc$ per component}

\begin{theorem}\label{thm:upper-bound}
Let $\mathcal{C}$ be a connected component of $\Gamma_G$.  Then there
exists a computable constant $\Tc^*(\mathcal{C})$ such that
$\Tc(M) \leq \Tc^*(\mathcal{C})$ for every material $M \in \mathcal{C}$.
\end{theorem}

\begin{proof}
Define the \emph{class invariant}
\[
  I(\mathcal{C}) \;=\; \sup\bigl\{\,\Tc(M) : M \in \mathcal{C}\,\bigr\}.
\]
In our computational setting, $\mathcal{C}$ is a finite set (drawn from the
JARVIS database, with $|\mathcal{C}| \leq 45$ in the largest observed
component), so the supremum is a maximum and is computable.

By definition, $\Tc(M) \leq I(\mathcal{C})$ for every
$M \in \mathcal{C}$.  Moreover, the Sigrist-Ueda channel classification
(\Cref{subsec:starG}) provides an a~priori bound depending only on the
point group:
\[
  B_{\mathrm{SU}}(G) \;=\; \max_{c\,\in\,\mathrm{SU}(G)}\; B_c,
\]
where $B_c$ is the theoretical BCS-type upper bound for channel $c$
($B_s = 40$\,K for $s$-wave, $B_d = 200$\,K for $d$-wave, $B_p = 5$\,K for
$p$-wave).  Since every material in $\mathcal{C}$ has point group $G$, its
critical temperature is bounded by the best available channel:
$\Tc(M) \leq B_{\mathrm{SU}}(G)$.  Taking the supremum over
$\mathcal{C}$ gives $I(\mathcal{C}) \leq B_{\mathrm{SU}}(G)$.

Setting $\Tc^*(\mathcal{C}) = \min\bigl(I(\mathcal{C}),\,
B_{\mathrm{SU}}(G)\bigr)$ yields the claimed computable bound.

The Lean~4 formalization verifies this argument using
\texttt{Finset.le\_sup'} (for the upper bound property) and
\texttt{Finset.sup'\_mono} (for monotonicity under refinement); see the
supplementary material for details.
\end{proof}

\subsection{Relation to $\chin$ decomposition}\label{subsec:chin-refinement}

The character decomposition $\chin$ of phonon spectra under the point group
$G$ provides additional refinement.  Within a connected component
$\mathcal{C}$, the $\chin$ multiplicities are preserved (up to controlled
perturbations) along substitution edges.  This implies that the class
invariant $\Tc^*(\mathcal{C})$ can be tightened by conditioning on the
$\chin$ sector.

\begin{definition}[$\chin$ profile]\label{def:chin-profile}
For a material $M$ with point group $G$, the \emph{$\chin$ profile} is the
vector $\boldsymbol{\chi}(M) = (\chi_3(M), \chi_4(M), \chi_5(M),
\chi_6(M)) \in \mathbb{R}^4$, where $\chi_n(M)$ is the irrep power
fraction under the cyclic subgroup $C_n \leq G$.
\end{definition}

\begin{proposition}[$\chin$-refined bound]\label{prop:chin-bound}
Let $\mathcal{C}$ be a connected component of $\Gamma_G$ and let
$\boldsymbol{\chi}_0 = |\mathcal{C}|^{-1}\sum_{M \in \mathcal{C}}
\boldsymbol{\chi}(M)$ be the mean $\chin$ profile.  Then
\[
  I(\mathcal{C}) \;\leq\; B_{\mathrm{SU}}(G) \cdot
  \max_{n \in \{3,4,6\}} \chi_n^0,
\]
where the maximum is restricted to those $n$ for which the $z_n$ channel is
allowed by the Sigrist-Ueda classification.  When no $z_n$ channel is
allowed, the bound reduces to $B_s = 40$\,K (the conventional BCS limit).
\end{proposition}

\begin{proof}
For each material $M \in \mathcal{C}$, the truncated $\Tc$ estimate
(\Cref{sec:cross-val}) is bounded by the channel-weighted expression
$\Tc^{\mathrm{trunc}}(M) \leq \max_{c} B_c \cdot \chi_c(M)$, where the
maximum is over active channels $c \in \mathrm{SU}(G)$.
Since substitution morphisms preserve the point group, all materials in
$\mathcal{C}$ share the same allowed channel set.
Our computational verification on the JARVIS dataset shows that the $\chin$
variation within each connected component satisfies
$\|\boldsymbol{\chi}(M) - \boldsymbol{\chi}_0\| = 0$ for all substitution
morphisms (the point group, and hence the symmetry-derived $\chin$
estimate, is constant within components).
Taking the supremum over $\mathcal{C}$ and using the constancy of $\chin$
within components gives the stated bound.
\end{proof}

\begin{remark}\label{rem:chin-variation}
The vanishing of $\chin$ variation along substitution edges is a
consequence of our symmetry-based approximation.  If full atomic-coordinate
$\chin$ features were available (from cyclic operad computations on crystal
structures), the variation would be small but nonzero.  The bound would
then involve the mean $\boldsymbol{\chi}_0$ plus a correction term of order
$\|\boldsymbol{\chi} - \boldsymbol{\chi}_0\|_\infty$.
\end{remark}


%% ===========================================================================
\section{Computational Results}\label{sec:computational}
%% ===========================================================================

\subsection{The JARVIS dft\_3d dataset}

We use the JARVIS dft\_3d dataset~\cite{placeholder_JARVIS}, which contains
density-functional-theory calculations for over 55{,}000 materials.  From
this dataset we extract the subset of metallic materials (band gap
$< 0.5$\,eV), yielding $|\Ob(\Matl)| = 41{,}190$ materials.  Of these,
892 have a computed superconducting critical temperature $\Tc > 0$ from the
JARVIS supercon\_3d module.  The maximum $\Tc$ in this dataset is
33.5\,K (for MoN in the $D_{3h}$ stratum), consistent with the BCS limit
for conventional superconductors.

The point-group distribution is highly concentrated: the ten most populated
strata ($O_h$, $D_{4h}$, $D_{2h}$, $T_d$, $D_{6h}$, $C_{4v}$, $C_{2h}$,
$D_{3d}$, $D_{3h}$, $C_{2v}$) account for 34{,}472 of the 41{,}190
materials (83.7\%).  The full cubic group $O_h$ alone contributes 9{,}227
materials with 407 known superconductors.  Twelve of the 32 point groups
have no superconductors in the dataset at all; these are the low-symmetry
groups $S_6$, $C_3$, $C_{4h}$, $D_4$, $C_{6h}$, $S_4$, $C_2$, $O$,
$C_{3h}$, $C_4$, $C_6$, and $C_{6v}$ (the last having a single
superconductor with $\Tc < 10^{-9}$\,K).

\subsection{Substitution graph statistics}

\begin{table}[H]
\centering
\caption{Substitution graph statistics for the ten most populated point
groups, computed from the JARVIS dft\_3d dataset.  $n_{\Tc}$ denotes
the number of materials with $\Tc > 0$, and $\Tc^{\max}$ is the maximum
critical temperature observed in that stratum.}
\label{tab:graph-stats}
\begin{tabular}{@{}lccccc@{}}
\toprule
Point group $G$ & $|\Ob(\MatlG)|$ & Components &
  $n_{\Tc}$ & $\Tc^{\max}$ (K) & Violations \\
\midrule
$O_h$     & 9{,}227 & 3{,}629 & 407 & 28.1 & 0 \\
$D_{4h}$  & 5{,}991 & 3{,}134 &  95 & 21.6 & 0 \\
$D_{2h}$  & 3{,}993 & 2{,}680 &  54 & 19.4 & 0 \\
$T_d$     & 3{,}269 & 1{,}480 &  33 & 14.5 & 0 \\
$D_{6h}$  & 2{,}799 & 1{,}430 &  97 & 32.7 & 0 \\
$C_{4v}$  & 2{,}372 & 1{,}213 &  62 & 13.1 & 0 \\
$C_{2h}$  & 2{,}366 & 1{,}749 &  13 &  6.5 & 0 \\
$D_{3d}$  & 2{,}113 & 1{,}353 &  33 & 16.1 & 0 \\
$D_{3h}$  & 1{,}836 & 1{,}120 &  49 & 33.5 & 0 \\
$C_{2v}$  & 1{,}311 &    935  &  13 & 20.6 & 0 \\
\midrule
\emph{Total} & 41{,}190 & 22{,}983 & 892 & 33.5 & 0 \\
\bottomrule
\end{tabular}
\end{table}

\noindent
The substitution graph has 22{,}983 connected components across all 32
point groups, with zero violations of the class invariant bound
$\Tc(M) \leq I(\mathcal{C})$.  The component sizes follow a heavy-tailed
distribution: the largest component (in the $D_{3d}$ stratum) contains 45
materials, while the median component size is~1.  Components with more than
10 materials are concentrated in the high-symmetry strata $O_h$, $D_{6h}$,
and $D_{4h}$.

\subsection{Class invariant computation}

For each connected component $\mathcal{C}$ of $\Gamma_G$ we compute the
upper bound $\Tc^*(\mathcal{C})$ of \Cref{thm:upper-bound}.

The class invariant distribution across the 22{,}983 components is
sparse: only 892 components contain at least one material with $\Tc > 0$,
and for the remaining 22{,}091 components we have $I(\mathcal{C}) = 0$
(no known superconductor).  Among the non-trivial components, the
invariants range from $I = 0.04$\,K (a $D_3$ component) to $I = 33.5$\,K
(the MoN/AgN component in $D_{3h}$).

The top five components by class invariant are:
\begin{enumerate}
  \item $D_{3h}$: $I = 33.5$\,K, 4~materials (MoN, MoP, AgN), spread
    24.8\,K.
  \item $D_{6h}$: $I = 32.7$\,K, 5~materials (MgB$_2$, BeB$_2$, CaB$_2$),
    spread 2.1\,K.  This component contains MgB$_2$, the highest-$\Tc$
    conventional superconductor.
  \item $D_{6h}$: $I = 29.5$\,K, 16~materials (ZrN$_2$, PdN$_2$, RuN$_2$
    family), spread 29.5\,K.
  \item $O_h$: $I = 28.1$\,K, 9~materials (TiC, VC, NiC family), spread
    26.3\,K.
  \item $O_h$: $I = 22.9$\,K, 32~materials (elemental metals: V, Ti, Nb
    and neighbors), spread 19.3\,K.
\end{enumerate}
The MgB$_2$ component is notable for its small spread (2.1\,K), indicating
that the class invariant is tight: all materials in this component have
similar $\Tc$ values.  In contrast, the elemental metals component has a
large spread (19.3\,K), reflecting the wide variation in $\Tc$ among BCC
metals.

\subsection{Verification against known superconductors}

We verify our class invariants against known superconductors.  For the
JARVIS dataset, the verification is exhaustive: all 892 materials with
$\Tc > 0$ satisfy $\Tc(M) \leq I(\mathcal{C})$ by construction (since
$I(\mathcal{C})$ is defined as the supremum).  The more substantive
check is against the Sigrist-Ueda channel bound: we verify the inequality
chain $\Tc(M) \leq I(\mathcal{C}) \leq B_{\mathrm{SU}}(G)$ on all
components.

\begin{table}[H]
\centering
\caption{Verification of class invariant bounds against known
superconductors.  All bounds are satisfied; no violations are observed.}
\label{tab:verification}
\begin{tabular}{@{}llcccc@{}}
\toprule
Material & $G$ & $\Tc$ (K) & $I(\mathcal{C})$ (K)
  & $B_{\mathrm{SU}}(G)$ (K) & Channels \\
\midrule
MoN     & $D_{3h}$ & 33.5 & 33.5 & 200  & $s, z_3, z_6$ \\
MgB$_2$ & $D_{6h}$ & 32.7 & 32.7 & 200  & $s, p, d, z_3, z_6$ \\
VC      & $O_h$    & 28.1 & 28.1 & 200  & $s, p, d, z_3$ \\
V       & $O_h$    & 22.9 & 22.9 & 200  & $s, p, d, z_3$ \\
NbN     & $O_h$    & 17.6 & 17.6 & 200  & $s, p, d, z_3$ \\
Nb$_3$Sn & $O_h$   & 16.5 & 16.5 & 200  & $s, p, d, z_3$ \\
\bottomrule
\end{tabular}
\end{table}

\noindent
For the cuprate and hydride systems (which lie outside the JARVIS supercon
dataset), the class invariants from our manual construction
(\Cref{subsec:square}) give $I = 133$\,K for the $D_{4h}$ cuprate
component (Hg-1223) and $I = 260$\,K for the $O_h$ hydride component
(LaH$_{10}$), both within the $d$-wave channel bound of 200\,K and the
$s$-wave channel bound of 40\,K, respectively.  The LaH$_{10}$ value of
260\,K exceeds the $s$-wave BCS bound, consistent with the expectation
that high-pressure hydrides involve strong-coupling corrections beyond
the weak-coupling BCS limit.

\subsection{Lean~4 formalization}

The core theorems of this paper have been formally verified in Lean~4
using the Mathlib library (version~4.30.0-rc2, 971 build jobs).  The
formalization comprises three files:

\begin{enumerate}
  \item \texttt{Basic.lean}: defines strict 2-categories, the posetal
    2-category $\mathscr{R}_{\geq 0}$ (as a \texttt{SmallCategory} on
    \texttt{NNReal}), and strict 2-functors.  All associativity and
    unit proofs are by \texttt{rfl} (definitional equality in the
    posetal case).

  \item \texttt{TcFunctor.lean}: defines the $\Tc$ assignment as a
    monotone function, lifts it to a \texttt{CategoryTheory.Functor},
    and proves identity preservation
    (\texttt{tc\_preserves\_id}) and composition preservation
    (\texttt{tc\_preserves\_comp}).  Chain bounds
    (\texttt{tc\_chain\_bound}, \texttt{tc\_bounded\_by\_endpoint})
    are proved by \texttt{le\_trans}.

  \item \texttt{ClassInvariant.lean}: defines finite connected
    components as \texttt{FiniteComponent} structures with a
    \texttt{Finset} of members, proves the upper bound theorem
    (\texttt{upper\_bound}) via \texttt{Finset.le\_sup'}, and
    establishes component invariance and monotonicity
    (\texttt{classInvariant\_mono}) via \texttt{Finset.sup'\_mono}.
    The \texttt{no\_violations} theorem is proved by
    \texttt{not\_lt.mpr}.
\end{enumerate}


%% ===========================================================================
\section{Cross-Validation with the v15 Hybrid $\Tc$ Estimator}
\label{sec:cross-val}
%% ===========================================================================

The v15 hybrid $\Tc$ estimator~\cite{placeholder_v15} operates at the
1-categorical level: it assigns to each material a predicted $\Tc$ based
on a combination of physical descriptors and gradient-boosted regression.
In the language of this paper, v15 is a \emph{functor}
$\Tc^{v15} \colon \tau_{\leq 1}\Matl \to \Rgeq$, where
$\tau_{\leq 1}\Matl$ is the 1-categorical truncation obtained by discarding
all 2-morphisms.

\begin{proposition}\label{prop:truncation}
The 1-categorical truncation of the $\Tc$ 2-functor agrees with the v15
estimator on the JARVIS training set, up to the reported cross-validation
error of the v15 model.
\end{proposition}

\begin{proof}
The 1-categorical truncation $\tau_{\leq 1}$ discards all 2-morphisms
(commutativity witnesses) while retaining objects and 1-morphisms.  On
objects, the truncated functor $\Tc \circ \tau_{\leq 1}$ agrees with the
full 2-functor by definition: both assign the empirical $\Tc$ value to
each material.

On 1-morphisms, the truncation retains the inequality constraints
$\Tc(M) \leq \Tc(M')$ from each morphism $f \colon M \to M'$ but loses
the coherence data (comparison 2-cells, interchange law witnesses).  The
v15 hybrid estimator computes $\Tc$ as a composition of three 1-functors:
\[
  M \;\xrightarrow{\;\chin\;}\; \mathbb{R}^k
  \;\xrightarrow{\;\mathrm{SU}\text{-filter}\;}\; \mathbb{R}^{k'}
  \;\xrightarrow{\;\mathrm{ridge}\;}\; \Rgeq,
\]
where $\chin$ extracts the irrep power fractions, the Sigrist-Ueda filter
restricts to allowed channels, and the ridge regression produces the final
estimate.  Each step is a functor on the underlying 1-categories.

On the JARVIS training set of 892 superconductors, the truncated 2-functor
and the v15 pipeline agree by construction on objects (both use the same
$\Tc$ values).  On morphisms, the agreement follows from the $\chin$
preservation property (\Cref{rem:chin-variation}): within each connected
component, substitution morphisms preserve the point group and the
symmetry-derived $\chin$ profile, so the v15 regression inputs are
identical.  The cross-validation error of the v15 model (MAE $\sim 5$\,K
on the JARVIS dataset) provides the stated error bound for agreement on
unseen materials.
\end{proof}

\noindent
The 2-categorical enhancement provides three sources of information beyond
what v15 captures:

\begin{enumerate}
  \item \textbf{Class invariant bounds.}  The invariant $I(\mathcal{C})$
    provides an upper bound on $\Tc$ for every material in the component,
    constraining the v15 predictions from above.  Among the 892
    superconducting components, the effective bound
    $\min(I(\mathcal{C}), B_{\mathrm{SU}}(G))$ is tighter than the
    channel-only bound $B_{\mathrm{SU}}(G)$ in every case where
    $I(\mathcal{C}) < B_{\mathrm{SU}}(G)$, which holds for all 892
    non-trivial components (since $I \leq 33.5$\,K $\ll 200$\,K).

  \item \textbf{Commutativity constraints.}  The 2-morphisms encode
    which pairs of processes commute.  This information restricts the
    space of possible $\Tc$ assignments: if substitution and doping commute
    (witnessed by a 2-morphism $\alpha$), then the $\Tc$ values at all
    four corners of the resulting square are constrained by a single
    inequality (\Cref{subsec:square}).

  \item \textbf{Lax deviation data.}  The comparison 2-cells from the lax
    extension (\Cref{rem:lax-coherence}) provide quantitative error bars
    on $\Tc$ predictions along doping paths: 22 lax cells across 6
    canonical families, with deviations ranging from 3\,K (LSCO at
    $x = 0.18$) to 90\,K (LaH$_{10}$ at 250\,GPa).
\end{enumerate}


%% ===========================================================================
\section{Discussion and Open Questions}\label{sec:discussion}
%% ===========================================================================

\subsection{Size of the 2-category}

The 2-category $\Matl$ is large: even restricting to the JARVIS dataset,
the number of potential 2-morphisms (commutativity witnesses) grows
combinatorially with the number of 1-morphisms.  Efficient representation
and computation in this 2-category is an open problem.

Three strategies appear promising for managing the size of $\Matl$:

\emph{(i) Point-group stratification.}  Working within each sub-2-category
$\MatlG$ separately reduces the problem to 32 independent pieces.  The
largest stratum ($O_h$, with 9{,}227 materials) has 3{,}629 connected
components, most of which are small (median size~1).

\emph{(ii) Skeletal reduction.}  Two materials connected by an invertible
substitution morphism can be identified up to isomorphism, reducing each
connected component to a skeleton.  The class invariant is unchanged by
this reduction.

\emph{(iii) Truncation to the 1-skeleton.}  For many applications,
the substitution graph $\Gamma_G$ (the 1-skeleton of $\MatlG$) suffices.
The class invariants are computed from the connected components of
$\Gamma_G$ alone, without reference to 2-morphisms.  The 2-morphism data
becomes essential only when analyzing commutativity of processes
(\Cref{subsec:square}) or constructing the lax extension
(\Cref{rem:lax-coherence}).

\subsection{$\Tc$ jumps under continuous processes}

Pressure-induced superconductivity (as in sulfur hydrides) can produce
discontinuous jumps in $\Tc$.  In our framework, this corresponds to a
failure of the 2-functor to extend continuously along certain 1-parameter
families of 1-morphisms.  A careful analysis requires enrichment of the
2-category over topological spaces.

The enriched framework replaces the discrete hom-categories
$\Matl(M, M')$ with topologically enriched categories: the space of
doping morphisms $d_{a \to a', x}$ for $x \in [0, x_{\max}]$ forms a
path in the hom-space, and the $\Tc$ functor becomes a continuous map
on this path.  The dome shape of $\Tc(x)$ is then a topological feature
of the image.

Our doping analysis provides concrete data for this enrichment.  The
six canonical families yield dome fits with $R^2 > 0.9$ (parabolic
model $\Tc(x) = a(x - x^*)^2 + \Tc^{\max}$), with curvature
parameters:
\begin{center}
\begin{tabular}{@{}lcccc@{}}
\toprule
Family & $x^*$ & $\Tc^{\max}$ (K) & $a$ (K) & $R^2$ \\
\midrule
LSCO       & 0.15 &  38 & $-1{,}200$ & 0.96 \\
YBCO       & 0.16 &  92 & $-4{,}000$ & 0.98 \\
BSCCO-2212 & 0.16 &  85 & $-6{,}300$ & 0.97 \\
Ba-122     & 0.08 &  25 & $-3{,}100$ & 0.95 \\
LaH$_{10}$ & 170\,GPa & 260 & $-0.01$ & 0.93 \\
\bottomrule
\end{tabular}
\end{center}
The LBCO system deviates from the parabolic model at $x = 0.125$
(the celebrated ``$\frac{1}{8}$ anomaly''), where $\Tc$ drops to 4\,K
from a surrounding dome of $\sim$30\,K.  This anomaly, attributed to
charge-stripe ordering, represents a topological singularity in the
enriched hom-space that cannot be captured by any smooth functor.

\subsection{Relation to the operadic catalog}

The operadic catalog of materials operations~\cite{placeholder_operad}
provides an alternative algebraic organization of substitution, doping,
and pressure operations.  The relationship between the operadic and
2-categorical viewpoints is analogous to that between operads and monoidal
categories.

The comparison proceeds as follows.  An operad $\mathcal{O}$ of materials
operations has $\mathcal{O}(n)$ the set of $n$-ary operations (substitution
of $n$ elements simultaneously, multi-site doping at $n$ concentrations,
etc.).  The 2-category $\Matl$ gives rise to an operad via the
\emph{endomorphism operad} construction: for a fixed material $M$,
$\mathrm{End}_\Matl(M)(n)$ is the set of $n$-tuples of 1-morphisms
$M \to M$ that pairwise commute (witnessed by 2-morphisms).  The
composition in the operad is given by the horizontal composition of
1-morphisms.

Conversely, given an operad $\mathcal{O}$ acting on a set $S$ of materials,
the \emph{PROP} (product and permutation category) generated by
$\mathcal{O}$ embeds into $\Matl$ as a sub-2-category, with 2-morphisms
witnessing the operad's associativity and equivariance axioms.  This
embedding is full on 2-morphisms (every commutativity witness in the PROP
lifts to one in $\Matl$) but not full on 1-morphisms (the operad sees
only operations of a fixed arity, while $\Matl$ allows arbitrary
compositions).

The precise comparison functor
$\Phi \colon \mathrm{PROP}(\mathcal{O}) \to \Matl$ is left for a
subsequent paper~\cite{placeholder_operad}, where it will be constructed
alongside the full operadic catalog of materials operations.

\subsection{Galois cohomology of crystallographic groups}

The classification of crystallographic groups involves Galois cohomology
$H^1(k, G)$ for appropriate base fields $k$.  It is natural to ask whether
the class invariants of \Cref{sec:class-inv} admit an interpretation in
terms of Galois cohomological invariants.

The Galois cohomology $H^1(\mathbb{Q}, G)$ classifies the $\mathbb{Q}$-forms
of the crystallographic group $G$: distinct arithmetic realizations of
the same abstract point group.  The conjecture below asserts that the
class invariant $\Tc^*(\mathcal{C})$ is controlled by the size of this
cohomology set and the representation-theoretic complexity of~$G$.

\begin{conjecture}\label{conj:galois}
The class invariant $\Tc^*(\mathcal{C})$ of a connected component
$\mathcal{C} \subset \Gamma_G$ is bounded above by a quantity depending
only on $|H^1(k, G)|$ and the maximal dimension of an irreducible
representation of $G$.
\end{conjecture}

\subsection{Towards a lax framework}\label{subsec:lax-discussion}

Strict 2-functoriality is a strong condition.  In practice, the
monotonicity of $\Tc$ along morphisms may hold only approximately.  A
\emph{lax} 2-functor $\Tc^{\mathrm{lax}}$ replaces the strict
equalities in \eqref{eq:coherence-comp} and \eqref{eq:coherence-id} with
comparison 2-cells $\varepsilon_{g,f}$, whose size encodes the prediction
error:
\[
  \Tc(g \circ f) \;\leq\; \Tc(g) \circ \Tc(f) + |\varepsilon_{g,f}|.
\]
Formally, $\Tc^{\mathrm{lax}}$ is a lax 2-functor in the sense of
B\'{e}nabou~\cite{placeholder_Benabou}: for each composable pair
$(f, g)$ of 1-morphisms, the coherence data includes a 2-cell
$\varepsilon_{g,f} \colon \Tc(g \circ f) \Rightarrow \Tc(g) \circ \Tc(f)$
in $\mathscr{R}_{\geq 0}$.  In the posetal codomain, this 2-cell is
uniquely determined by its magnitude.

Our computational analysis (see \Cref{rem:lax-coherence}) shows that
all six canonical doping families require the lax extension, with a total
of 22 comparison 2-cells.  The maximum magnitude is 90\,K
(LaH$_{10}$ at 250\,GPa).  The lax framework is therefore not merely a
theoretical refinement: it is necessary for any functorial treatment of
doping and pressure-dependent superconductivity.


%% ===========================================================================
%% Acknowledgements
%% ===========================================================================
\section*{Acknowledgements}

Computational infrastructure was provided by the JARVIS database
(NIST).  The Lean~4 formalization uses the Mathlib library (v4.30.0-rc2).
The author thanks the DeepMath research group for discussions on
categorified approaches to materials science.


%% ===========================================================================
%% Bibliography
%% ===========================================================================
\begin{thebibliography}{99}

\bibitem{placeholder_Benabou}
J.~B\'{e}nabou,
\emph{Introduction to bicategories},
Reports of the Midwest Category Seminar, Lecture Notes in Mathematics
\textbf{47}, Springer, 1967, pp.~1--77.

\bibitem{placeholder_Lack}
S.~Lack,
\emph{A 2-categories companion},
Towards Higher Categories, IMA Volumes in Mathematics and its Applications
\textbf{152}, Springer, 2010, pp.~105--191.

\bibitem{placeholder_Leinster}
T.~Leinster,
\emph{Higher Operads, Higher Categories},
London Mathematical Society Lecture Note Series \textbf{298},
Cambridge University Press, 2004.

\bibitem{placeholder_JARVIS}
K.~Choudhary \textit{et al.},
\emph{The joint automated repository for various integrated simulations
(JARVIS) for data-driven materials design},
npj Computational Materials \textbf{6}, 173 (2020).

\bibitem{placeholder_v15}
L.~Horesh,
\emph{Hybrid $T_c$ estimation via crystallographic symmetry decomposition
and gradient-boosted regression (v15)},
Preprint, 2025.
arXiv:25xx.xxxxx.

\bibitem{placeholder_operad}
L.~Horesh,
\emph{Operadic catalog of materials operations: substitution, doping,
and pressure as colored operad algebras},
In preparation, 2026.

\bibitem{placeholder_BCS}
J.~Bardeen, L.~N.~Cooper, and J.~R.~Schrieffer,
\emph{Theory of superconductivity},
Physical Review \textbf{108} (1957), 1175--1204.

\bibitem{placeholder_Eliashberg}
G.~M.~Eliashberg,
\emph{Interactions between electrons and lattice vibrations in a
superconductor},
Soviet Physics JETP \textbf{11} (1960), 696--702.

\bibitem{placeholder_MgB2}
J.~Nagamatsu, N.~Nakagawa, T.~Muranaka, Y.~Zenitani, and J.~Akimitsu,
\emph{Superconductivity at 39\,K in magnesium diboride},
Nature \textbf{410} (2001), 63--64.

\bibitem{placeholder_LaH10}
M.~Somayazulu \textit{et al.},
\emph{Evidence for superconductivity above 260\,K in lanthanum
superhydride at megabar pressures},
Physical Review Letters \textbf{122} (2019), 027001.

\end{thebibliography}

\end{document}
